% Section 3. Methods. From single labels to repeated labels.
% Goal. Each slide answers the question the previous slide raises. See STORY.md.

\section{Methods}

% Slide 3.1. Source: Methods 4.2.1 (sec:edl-model-setup, eq:edl-dirichlet), 4.3.1
% (sec:ebnn-model, eq:ebnn-hierarchy, eq:edl-bnn-predictive), 4.3.2
% (sec:ebnn-likelihood-equivalence, eq:ebnn-likelihood-equivalence), 4.3.5
% (sec:ebnn-uncertainty-decomposition).
\begin{frame}{The Categorical Evidential Bayesian Neural Network}
	\textbf{Message.} The Cat-eBNN is EDL's Dirichlet network with a posterior over the shared network weights. It represents all three terms.

	\medskip
	\textbf{TODO: content.} Model sketch as a small graph. Likelihood is the Dirichlet mean $m_y(x;w)$.
\end{frame}

% Slide 3.2. Source: Methods 4.3.3 (sec:ebnn-identifiability, eq:ebnn-likelihood-invariance),
% 4.3.4 (sec:gamma-prior, eq:gamma-prior-weights, eq:assumed-total-concentration).
% Background 2.10 (sec:bg-gamma, eq:gamma-mode).
\begin{frame}{Single Labels Cannot Inform the Split. The Gamma Prior}
	\textbf{Message.} The single label likelihood is invariant to the total concentration parameter at a fixed mean vector. So we state the assumption with a Gamma prior on $\alpha_0(x;w)$.

	\medskip
	\textbf{TODO: content.} Invariance in one line. Gamma prior with prior mode and prior standard deviation. Honest scope sentence.
\end{frame}

% Slide 3.3. Source: Background 2.5 (sec:bg-count-models, eq:count-vector). Methods 4.4.1
% (sec:mbnn-model, eq:mbnn-observation), 4.4.2 (sec:dm-ebnn-model, eq:dm-ebnn-observation),
% 4.4.5 (sec:count-uncertainty-decomposition).
\begin{frame}{Repeated Labels. The Dirichlet Multinomial eBNN}
	\textbf{Message.} Repeated labels, aggregated by class into count vectors, are a different observation. The DM-eBNN models them with a Dirichlet multinomial likelihood. The M-BNN is the Bayesian baseline with a multinomial likelihood.

	\medskip
	\textbf{TODO: content.} Count vector $n$ and count total $R$. The two observation models side by side. Spell out M-BNN.
\end{frame}

% Slide 3.4. Source: Background 2.5.2 (eq:dm-moments). Methods 4.4.3
% (sec:dm-ebnn-concentration-information, eq:dm-ebnn-posterior-odds), 4.4.4
% (sec:count-comparison, eq:count-variance-ratio).
\begin{frame}{Count Data Can Inform the Split}
	\textbf{Message.} At a fixed mean vector and a count total above one, the count variance depends on the total concentration parameter. So the count likelihood can inform it. The M-BNN has a fixed count variance.

	\medskip
	\textbf{TODO: content.} Dirichlet multinomial variance in one line. Variance ratio in one line. The Gamma prior stays part of the posterior.
\end{frame}
